\chapter{宇宙学}\label{chpt:cosmos}
空间均匀的宇宙学模型

宇宙学的早期发展毫无疑问是 Einstein 带头。
观测结果建议，宇宙在大尺度上均匀且各向同性。这使得不能用非零质量密度和常数场强求解 Poisson 方程，因为试图计算某点场强时，要涉及对充斥全空间的物质做积分，这是反常积分，需用有限域逼近。简单计算就可发现结果取决于趋于无穷的方式（“为零”或“发散”），则称为 Neumann-Zeeliger 疑难，说明积分不良定义。可见 Newton 引力论很难发展出宇宙学。
相较而言，广义相对论可以自然地研究均匀各向同性的解及其扰动。
比如，\textbf{FLRW 解}就是膨胀宇宙的一阶近似。

1967 年 Hawking 证明了膨胀宇宙在时间上回溯可到达奇点，这被称为\textbf{宇宙大爆炸}(the big bang)。
\section{运动学}
\section{动力学}
\section{热力学}
\section{标准模型疑难}
\section{暗能量}
\section{暴涨理论}